\chapter*{Abstract}

\small
\noindent LIMA, T. A. S. \emph{Comparative analysis between monolithic architecture and
microservices architecture}. 2026. Monografia (Graduação em Ciência da Computação)
--- Centro Universitário de Formiga -- UNIFOR-MG, Formiga, 2026.

\normalsize \vspace{1cm} \onehalfspacing

\begin{otherlanguage}{english}
\noindent Choosing between the monolithic and the microservices architectures is one of the
most relevant decisions in Web application development, and the studies that compare both
approaches report divergent results. This work compared the two architectures through two
functionally equivalent versions of an e-commerce application, developed in Java with Spring
Boot, featuring authentication, product catalog, orders, payments integrated with Mercado Pago,
file storage and asynchronous communication with Kafka. Both versions were executed in Docker
containers with the same resource budget and evaluated regarding performance, scalability,
maintainability, development complexity and deployment, using load and stress tests generated
by k6 and repeated five times. The monolithic version showed lower response times in all
scenarios, sustained 200 iterations per second in the mixed scenario, against 150 for the
microservices, consumed less CPU and about four times less memory and started 4.5 times
faster. With two instances behind a load balancer, the monolith doubled its throughput. The
microservices version provided fault isolation between the catalog and the purchase flow and
allowed scaling only the most demanded services, but required about 50\% more code and its
capacity was limited by the share of resources assigned to each service. It is concluded that, for small to
medium applications developed by small teams with limited resources, the monolithic
architecture is more advantageous, while the benefits of microservices depend on a context with
multiple teams and an infrastructure able to scale each service independently.

\vfill%
\noindent Keywords: Software architecture. Microservices. Monolithic architecture. E-commerce.
\end{otherlanguage}
